\clearpage\begingroup\HistoricalHead{227}{R}{BESTIMMUNG DER BAHNEN DER HIMMELSKÖRPER.}
\begin{center}[5.]\end{center}
\noindent\hspace*{\parindent}Mit den solchergestalt gefundenen Elementen, die den beiden äussern Beobachtungen genau Genüge leisten, berechnet man nun den geocentrischen Ort für die Zeit der mittlern Beobachtung. Stimmt derselbe ganz mit dem beobachteten, so hat man die ganze Bahn schon richtig; wo nicht, so muss man die Elemente durch neue Hypothesen verbessern. — Übrigens [wird] man
\fancyfoot[R]{\small 29\textsuperscript{*}}
\par\clearpage
\HistoricalHead{228}{L}{THEORETISCHE ASTRONOMIE. NACHLASS.}\SourceLeaf{233}
\noindent zur Berechnung des geocentrischen Orts aus dem heliocentrischen in der Bahn folgende Formeln sehr brauchbar finden.
\[
\begin{aligned}
r\cos(v-\AstNode)+R\cos(V-\AstNode)&=\Delta\cos(\alpha-\AstNode),\\
r\sin(v-\AstNode)\mathbin{.}\cos i+R\sin(V-\AstNode)&=\Delta\sin(\alpha-\AstNode),\\
r\sin(v-\AstNode)\mathbin{.}\sin i&=\Delta\mathbin{.}\Theta,
\end{aligned}
\]
\[
\begin{aligned}
&\Delta\surd(1+\Theta\Theta)={}\\
&\surd\Bigl\{rr+RR+2rR\bigl(\cos(v-\AstNode)\cos(V-\AstNode)\\
&\hspace{3em}+\sin(v-\AstNode)\sin(V-\AstNode)\cos i\bigr)\Bigr\}={}\\
&\surd\Bigl\{(r+R)^2\bigl(\cos\tfrac12i^2\cos\tfrac12(V-v)^2\\
&\hspace{3em}+\sin\tfrac12i^2\cos\bigl(\tfrac12(V+v)-\AstNode\bigr)^2\bigr)\\
&\hspace{2em}+(r-R)^2\bigl(\cos\tfrac12i^2\mathbin{.}\sin\tfrac12(V-v)^2\\
&\hspace{3em}+\sin\tfrac12i^2\sin\bigl(\tfrac12(V+v)-\AstNode\bigr)^2\bigr)\Bigr\}.
\end{aligned}
\]
\begin{center}
[6.]\par\medskip
\textls[100]{Erste Verbesserungsmethode.}
\end{center}
\noindent{\small [S. 6]}\par
\noindent\hspace*{\parindent}Man mache zwei neue Hypothesen, indem man die curtirten Abstände von der Erde für die äussern Beobachtungen etwas verändert; berechne daraus die Elemente der Bahn, und hieraus die mittlere Beobachtung. Durch Vergleichung der für diese nach den drei Hypothesen gefundenen Fehler wird man sodann die Abstände viel genauer ableiten und daraus der Wahrheit viel nähere Elemente finden. Da diese Rechnungen den vorhin beschriebenen ganz ähnlich sind, so ist darüber nichts weiter hinzuzusetzen.
\begin{center}
[7.]\par\medskip
\textls[100]{Zweite Verbesserungsmethode.}
\end{center}
\noindent\hspace*{\parindent}Man verändere die aus der ersten Hypothese abgeleitete Neigung der Bahn und den Ort des Knoten etwas; berechne damit die zwei heliocentrischen Örter in der Bahn und Abstände von der Sonne; aus den beiden äussern und der Zwischenzeit berechne man die Elemente und den heliocentrischen Ort und Abstand der mittlern Beobachtung. Dasselbe berechne man nach den drei Hypothesen aus dem beobachteten geocentrischen Orte, und leite aus der Vergleichung der Fehler die nähern Werthe für Neigung und Knoten ab. Hier müssen wir also auflösen folgende
\clearpage
\HistoricalHead{229}{R}{BESTIMMUNG DER BAHNEN DER HIMMELSKÖRPER.}\SourceLeaf{234}
\begin{center}\textls[100]{Aufgabe.}\end{center}
\noindent\hspace*{\parindent}Aus dem geocentrischen Orte, Neigung der Bahn und Länge des Knoten, den heliocentrischen Ort in der Bahn und die Abstände von der Sonne und Erde zu finden\textsuperscript{*})
\begingroup\renewcommand{\thefootnote}{}\footnotetext{\textsuperscript{*}) Der letztere ist eigentlich bloss zur genauen Bestimmung der Aberration nöthig.}\endgroup
\[
\frac R\Delta=\frac{\sin(\alpha-\AstNode)-\dfrac\Theta g}{\sin(V-\AstNode)},
\]
\[
\frac r\Delta=\surd\left(1+\Theta\Theta+\frac{RR}{\Delta\Delta}-2\frac R\Delta\cos(V-\alpha)\right),
\]
\[
\sin(v-\AstNode)=\frac\Theta{\dfrac r\Delta\sin i},
\]
\[
\begin{aligned}
\cos(v-\AstNode)&=\frac\Delta r\cos(\alpha-\AstNode)-\frac Rr\cos(V-\AstNode)\\
&=\frac\Delta r\left(\frac{\sin(V-\alpha)}{\sin(V-\AstNode)}+\frac\Theta g\operatorname{cotg}(V-\AstNode)\right),
\end{aligned}
\]
\[
\begin{aligned}
\operatorname{cotg}(v-\AstNode)&=\frac{\sin i}\Theta\frac{\sin(V-\alpha)}{\sin(V-\AstNode)}+\cos i\operatorname{cotg}(V-\AstNode)\\
&=\frac{\cos i\mathbin{.}g}{\Theta\sin(V-\AstNode)}\left(\sin(V-\alpha)+\frac\Theta g\cos(V-\AstNode)\right).
\end{aligned}
\]
\clearpage\endgroup
